\documentclass[10pt,a4paper]{amsart}
\usepackage[margin=19mm]{geometry}
\usepackage{amsmath,amssymb,mathtools}
\usepackage[scr=rsfs,cal=euler]{mathalfa}
\usepackage{xfrac}
\usepackage[unicode,hidelinks,bookmarks=false]{hyperref}
\newcommand{\ie}{i.e.\ }
\newcommand{\cf}{cf.\ }
\newcommand{\resp}{respectively\ }
\newcommand{\Subsection}{\subsection}
\providecommand{\nobreakdash}{\nobreak\hbox{-}}
\setlength{\emergencystretch}{2em}
\multlinegap0pt
\usepackage{booktabs}
\usepackage{subfig}
\usepackage{mathtools}
\usepackage{amssymb}
\usepackage[shortlabels]{enumitem}
\usepackage{bm}
\usepackage{nicefrac}
\newtheorem{prop}{Proposition}
\newcommand{\C}{\mathbb{C}}
\newcommand{\R}{\mathbb{R}}
\newcommand{\Tr}{\operatorname{Tr}}
\newcommand{\im}{\operatorname{Im}}
\newcommand{\re}{\operatorname{Re}}
\title{\large Kudla Millson lift of toric cycles and restriction of Hilbert modular forms.}
\author{Romain Branchereau}
\date{Source: arXiv:2401.16996v4 (CC BY 4.0)}
\begin{document}
\maketitle

\noindent\emph{Source and licence.} \large Kudla Millson lift of toric cycles and restriction of Hilbert modular forms.. Version arXiv:2401.16996v4 (CC BY 4.0), distributed by arXiv under the Creative Commons Attribution 4.0 International licence, http://creativecommons.org/licenses/by/4.0/, as recorded in the arXiv licence metadata for this exact version and shown on the abstract page https://arxiv.org/abs/2401.16996v4. Full licence notice: \texttt{licenses/arxiv-2401.16996-CC-BY-4.0.txt}.\par\smallskip

\noindent\textit{Editorial scope.} Counted unit: the article's prop sign1 (source0.tex lines 683--684). The article's complete original proof of it is delivered with it (source0.tex lines 685--716, one \texttt{proof} environment of 32 lines). No mathematics is written, repaired or re-derived: the statement and the proof are the article's own lines, in the article's own notation, and no retained line is substituted or re-flowed.  No further source-local context is needed: every cross-reference of every retained line resolves inside the two delivered spans.  Nothing was omitted from any retained span, so the package carries no omission record.  No retained line contains a citation, so the wrapper carries no bibliography and no bibliography entry.  No retained line contains a figure environment, an image-inclusion command or drawing code, so no image is delivered and the package declares no asset dependency.  Editorial formatting only, in addition: an AMS \texttt{amsart} preamble that restates the article's own declarations and macros used by the retained lines, namely the environments \texttt{prop} on separate counters, together with the standard packages those lines need.  The retained environments print in this document as Proposition 1 in order of appearance, whereas the article prints the same items with the numbering of its own full document.  The complete native source of the article as distributed in the arXiv e-print for this exact version is bundled unchanged as \texttt{sources/153744/source0.tex}.
\begin{prop} \label{sign1} The signature of $E_\R$ is $(n_1+2r_\alpha+2n_3,n_1+2s_\alpha+2n_3)$. 
\end{prop}

\begin{proof}
We have an isomorphism of quadratic spaces $(E_\R, Q_\alpha) \simeq \bigoplus (E_\sigma,Q_{\sigma(\alpha)})$. We only have to find the signature at each place $\sigma$. 
\begin{enumerate}
    \item If $\sigma \in S_1$, then $E_\sigma=\R \times \R$ and the involution permutes the two factors. Hence 
    \begin{align}
        Q_{\sigma(\alpha)}\left ( \begin{pmatrix}
            t_1 \\ t_1'
        \end{pmatrix},\begin{pmatrix}
            t_2 \\ t_2'
        \end{pmatrix}\right )=\sigma(\alpha)(t_1t_2'+t_1't_2)
    \end{align}
    and $E_\sigma$ has signature $(1,1)$.
    \item If $\sigma \in S_2$, then $E_\sigma=\C \simeq \R^2$ and the involution is the complex conjugation. Hence 
    \begin{align}
        Q_{\sigma(\alpha)}\left ( z_1,z_2\right )=\sigma(\alpha)\Tr_{\C/\R}(z_1\bar{z_2})=\sigma(\alpha)\left [\re(z_1)\re(z_2)+\im(z_1)\im(z_2) \right ],
    \end{align}
    and the signature of $E_\sigma$ is $(2,0)$ if $\sigma(\alpha)>0$ and $(0,2)$ if $\sigma(\alpha)<0$.
    \item If $\sigma \in S_3$, then $E_\sigma=\C\times \C$ and the involution permutes the two factors. Hence 
    \begin{align}
        Q_{\sigma(\alpha)}\left ( \begin{pmatrix}
            z_1 \\ w_1
        \end{pmatrix},\begin{pmatrix}
            z_2 \\ w_2
        \end{pmatrix}\right )=\sigma(\alpha)(z_1w_2+z_2w_1)
    \end{align}
    and $E_\sigma$ has signature $(2,2)$. 
\end{enumerate}
It follows that the signature of $E_\R$ is
\begin{align}
    n_1(1,1)+r_\alpha(2,0)+s_\alpha(0,2)+n_3(2,2)=(n_1+2r_\alpha+2n_3,n_1+2s_\alpha+2n_3)
\end{align}
\end{proof}

\end{document}
